\documentclass[11pt,a4paper]{article}

% ---------------------------------------------------------------------------
% Glossary of Acronyms & Symbols
% Copper Alloy Digital Twin -- uncertainty-aware surrogate pipeline
% ---------------------------------------------------------------------------

\usepackage[utf8]{inputenc}
\usepackage[T1]{fontenc}
\usepackage{lmodern}          % scalable fonts: required before microtype (MiKTeX)
\usepackage[margin=1in]{geometry}
\usepackage{microtype}
\usepackage{amsmath}
\usepackage{amssymb}
\usepackage{xcolor}
\usepackage{enumitem}
\usepackage{titlesec}
\usepackage{booktabs}
\usepackage[colorlinks=true,linkcolor=blue!55!black,urlcolor=blue!55!black]{hyperref}

\newcommand{\code}[1]{\texttt{\small #1}}

% term list: bold term, hanging indent
\newenvironment{terms}
  {\begin{description}[style=unboxed,leftmargin=1.2em,itemsep=0.45em,parsep=0pt,
                       font=\normalfont\bfseries]}
  {\end{description}}

\titleformat{\section}{\large\bfseries}{\thesection}{0.6em}{}
\titlespacing*{\section}{0pt}{1.4em}{0.6em}

\title{\textbf{Copper Alloy Digital Twin}\\[0.3em]
\large Glossary of Acronyms and Symbols}
\author{Uncertainty-aware surrogate pipeline}
\date{\today}

\begin{document}
\maketitle

\noindent
Quick reference for the acronyms, symbols, and column names used across the
notebooks, modules, and exported CSVs. Related concepts are grouped together
(e.g.\ $\kappa$ and LCB). For the mathematical derivations see
\code{copper\_digital\_twin\_v4}; for the module-level walkthrough see
\code{code\_architecture}.

\vspace{0.5em}
\hrule

% ===========================================================================
\section{Architecture}
% ===========================================================================

\begin{terms}
  \item[MBC --- Model-Based Controller] The decision layer that chooses process
    parameters ($T$, $t$ for annealing; the sequence of passes for cold
    drawing) from the surrogates' predictive distributions.

  \item[Surrogate] A machine learning model emulating an expensive physical
    process. Nine in total (three properties $\times$ three processes).

  \item[Model A / Model B] Model A predicts the mean $\hat\mu$ and the
    epistemic variance; Model B predicts the aleatoric variance. Together they
    form the predictive distribution.

  \item[Bundle] The unit of persistence: \code{artifacts.pkl} $+$
    \code{model\_a/} $+$ \code{model\_b/}. The only interface between training
    and the controllers.
\end{terms}

% ===========================================================================
\section{Ensemble and weights}
% ===========================================================================

\begin{terms}
  \item[NNLS --- Non-Negative Least Squares] Recovers the AutoGluon
    WeightedEnsemble weights by solving $\min_{w \ge 0}\lVert Aw - b\rVert$,
    with no dependence on AutoGluon's internal API.

  \item[$w$ (weights)] Convex weights over the base learners ($w_m \ge 0$,
    $\sum_m w_m = 1$). They act as an empirical posterior over models.

  \item[\code{use\_weighted\_variance}] Flag toggling between $w$-weighted
    statistics (default) and uniform statistics (robustness check).
\end{terms}

% ===========================================================================
\section{Uncertainty decomposition}
% ===========================================================================

\begin{terms}
  \item[$\sigma^2_{\text{epist}}$ (epistemic)] Disagreement among the base
    learners. Model ignorance --- \emph{reducible} with more data or a better
    model class.

  \item[$\sigma^2_{\text{aleat}}$ (aleatoric)] Intrinsic process noise ---
    \emph{irreducible} with more data.

  \item[$\sigma_{\text{total}}$] $\sqrt{\sigma^2_{\text{epist}} +
    \sigma^2_{\text{aleat}}}$, after recalibration. The standard deviation of
    the predictive distribution --- \emph{not} the error of that particular
    prediction.

  \item[$\tilde r^2$] Model B's training target,
    $\max\!\big((y - \hat\mu_{\text{oof}})^2 - \hat\sigma^2_{\text{epist}},\,0\big)$.
    Subtracting the epistemic part from the squared residual isolates the
    aleatoric one.

  \item[OOF --- Out-Of-Fold] A prediction made by a model that did \emph{not}
    see that row during training. Honest estimate of generalisation error;
    the basis of NNLS, Model B, and calibration.
\end{terms}

% ===========================================================================
\section{Calibration}
% ===========================================================================

\begin{terms}
  \item[$\alpha$ (alpha)] Nominal coverage promised by an interval
    ($0.5$, $0.8$, $0.9$, $0.95$): ``90\% of true values should fall inside''.

  \item[coverage] Empirical coverage --- the fraction that actually fell
    inside. Calibrated when $\widehat{\text{coverage}} \approx \alpha$.

  \item[$c$ (\code{recalibration\_c})] Single scalar correcting
    $\sigma \mapsto c\,\sigma$ so empirical coverage matches nominal at
    $\alpha = 0.9$. $c > 1$ means the model was overconfident.

  \item[$z_\alpha$] Standard normal quantile ($z_{0.9} = 1.645$,
    $z_{0.95} = 1.960$). Sets the interval half-width.
\end{terms}

% ===========================================================================
\section{Physical constraints}
% ===========================================================================

\begin{terms}
  \item[\code{y\_min} / \code{y\_max}] Physical support of the target
    (e.g.\ the $\sim$105 ceiling on \%IACS; grain size $> 0$).

  \item[Truncated Gaussian] The predictive Gaussian clipped to
    $[y_{\min}, y_{\max}]$, applied \emph{at inference only} (training
    untouched).

  \item[\code{mu\_trunc} / \code{sigma\_trunc}] Moments of that truncated law.
    Always $\sigma_{\text{trunc}} < \sigma_{\text{total}}$ --- truncation
    narrows.

  \item[\code{p\_above\_max}] Diagnostic: the probability mass the untruncated
    law placed \emph{beyond} the ceiling. High values mean the prediction is
    leaning on the wall --- treat with suspicion.
\end{terms}

% ===========================================================================
\section{Out-of-domain handling (OOD)}
% ===========================================================================

\begin{terms}
  \item[OOD --- Out-Of-Distribution] An input outside the training cloud, where
    tree ensembles extrapolate flat and overconfidently.

  \item[Mahalanobis distance $d$] Distance accounting for the covariance of
    the features; measures how unusual an input is.

  \item[$d_{\text{ref}}$] The $95^{\text{th}}$ percentile of the training
    distances --- the edge of the cloud.

  \item[$\gamma$ (gamma) and $m(x)$]
    $m(x) = \min\!\big(1 + \gamma \cdot \text{relative excess},\,
    m_{\max}\big)$, the multiplier applied to $\sigma_{\text{epist}}$. Exactly
    $1$ inside the cloud; grows smoothly outside, capped at $3$.
\end{terms}

% ===========================================================================
\section{Decision layer (what the controller consumes)}
% ===========================================================================

\begin{terms}
  \item[$\delta$ (delta)] Risk tolerance. \textbf{Filters}: keeps only cells or
    routes with $\Pr(\text{success}) \ge 1 - \delta$. Binary (pass / fail).

  \item[\code{pr\_success\_all}] Probability of meeting all setpoints. On the
    annealing grid it is a product (independence assumed); for cold-drawing
    routes it is the empirical \emph{joint} probability over the shared Monte
    Carlo trajectories --- no independence assumption.

  \item[Chance constraint] The form $P(y \ge y_{\text{target}}) \ge 1-\delta$,
    replacing the deterministic controller's binary
    \code{pred >= min} test.

  \item[$\kappa$ (kappa) and LCB --- Lower Confidence Bound]
    Always paired: $\mathrm{LCB} = \hat\mu - \kappa\,\hat\sigma$. With
    $\kappa = 1.645$ there is a $\sim$95\% chance the true value lies above it.
    \textbf{Ranks, never filters} --- among equal means it prefers the smaller
    $\sigma$. It is a linear approximation of what \code{pr\_success\_all}
    already does exactly.
\end{terms}

% ===========================================================================
\section{Chained propagation}
% ===========================================================================

\begin{terms}
  \item[Upstream / downstream] Producer vs.\ consumer. Grain size is upstream
    of IACS; pass $i$ is upstream of pass $i+1$.

  \item[MC (Monte Carlo) / $K$] Drawing $K = 200$ trajectories from the
    upstream distribution instead of propagating only its mean. Without it,
    the inherited variance is silently dropped.

  \item[\code{sigma2\_propagated}] The share of the downstream $\sigma^2$
    contributed by upstream uncertainty. Its fraction of the total answers:
    ``would more grain-size data tighten the IACS forecast?''
\end{terms}

% ===========================================================================
\section{Cold drawing}
% ===========================================================================

\begin{terms}
  \item[One-step ahead] Every pass predicted from \emph{true} inputs. Per-pass
    model quality, with no compounding.

  \item[Multistep rollout] The prediction of pass $i$ feeds pass $i+1$. The gap
    against one-step isolates error \textbf{compounding}.

  \item[\code{compounding\_ratio}] How much the final error grows relative to
    what independent errors would give. $\approx 1.0$ means $\sqrt{n}$-like
    growth; $> 1.0$ means correlated errors compounding faster.

  \item[GroupKFold / group folds] Whole wires per fold. Rows of the same wire
    are dependent; random folds leak information and corrupt even Model B's
    targets.

  \item[DAG / DFS] The directed acyclic graph of reachable diameters and the
    depth-first search enumerating routes $D_0 \rightarrow D_f$.

  \item[\code{rr\_max}] Maximum admissible reduction ratio per pass (\%).
    Prunes infeasible edges during graph construction.

  \item[\code{total\_strain} / \code{reduction\_ratio}] Derived per-pass
    geometry: $2\ln(D_i/D_f)$ and $(1 - D_f^2/D_i^2)\times 100$. Always
    recomputed from the measured diameters.
\end{terms}

% ===========================================================================
\section{Classical metrics}
% ===========================================================================

\begin{terms}
  \item[MAE / RMSE / MAPE] Mean absolute error, root mean squared error, mean
    absolute percentage error. Still reported alongside calibration --- they
    answer different questions. Expected consistency:
    $\mathbb{E}\big[\sigma_{\text{total}}\big] \approx \text{RMSE}$.
\end{terms}

\end{document}